% Part: sets-functions-relations
% Chapter: infinite
% Section: dedekind-induction

\documentclass[../../../include/open-logic-section]{subfiles}

\begin{document}

\olfileid{sfr}{infinite}{induction}
\olsection[حسابی استقرا]{ڈیڈیکائنڈ الجبرے تے حسابی استقرا}

ہُن بنیادی گل ایہ اے کہ اک ڈیڈیکائنڈ الجبرا---سچ پچھو تاں
\emph{ہر} ڈیڈیکائنڈ الجبرا---قدرتی عدداں دا قائم مقام بن سکدا اے۔
ایہ اگلے فوری نتیجے دی بدولت اے:

\begin{thm}[حسابی استقرا]\ollabel{thm:dedinfiniteinduction}
$N, s, o$ رل کے اک ڈیڈیکائنڈ الجبرا بناؤن۔ فیر ہر سیٹ $X$ لئی:
	\begin{center}
		جے $o \in X$ تے $(\forall {n} \in N \cap X){s}({n}) \in X$، {تاں} $N \subseteq X$۔
	\end{center}
\end{thm}

\begin{proof}
ڈیڈیکائنڈ الجبرے دی تعریف مطابق $N = \closureofunder{s}{o}$۔ ہُن جے
${o} \in X$ تے $(\forall {n} \in N)(n \in X \lif
{s}({n}) \in X)$ دونیں ہون، تاں حامل تے دِتے سیٹ دا اشتراک، خاص رکن
نوں رکھن والا جانشین-بند حامل ذیلی سیٹ اے؛ ایس لئی
$N = \closureofunder{s}{o} \subseteq X$۔
% PNBINF-003 ماخذ دی درستی: قضیہ X نوں ہر سیٹ مندا اے، سو بندش دا چھوٹاپن N نال X دے اشتراک اُتے لگدا اے؛ ماخذ دا ثبوت X نوں سدھا s-بند آکھ کے حامل توں باہر دے ممکنہ رکناں نوں نظر انداز کردا اے۔
\end{proof}

کیوں جو استقرا قدرتی عدداں دی پہچان اے، نکتہ ایہ اے: کوئی وی
ڈیڈیکائنڈ لامتناہی سیٹ دِتا ہووے، اسی اک ڈیڈیکائنڈ الجبرا بنا کے اوس
الجبرے نوں قدرتی عدداں دے قائم مقام دے طور تے ورت سکدے آں۔

منیا کہ \olref{thm:dedinfiniteinduction} استقرا نوں \emph{سیٹ تھیوری}
دے لفظاں وچ بیان کردا اے۔ پر اسی اصول نوں اوہناں لفظاں وچ سوکھ نال
لکھ سکدے آں جیہڑے شاید ہور جانے پچھانے ہون:

\begin{cor}\ollabel{natinductionschema}
$N, s, o$ رل کے اک ڈیڈیکائنڈ الجبرا بناؤن۔ فیر ہر فارمولے
$\phi(x)$ لئی، جیہدے وچ پیرامیٹر ہو سکدے نیں:
\begin{center}
	جے $\phi(o)$ تے $(\forall {n} \in N)(\phi(n)\lif
	\phi({s}({n})))$، {تاں} $(\forall n \in N)\phi(n)$
\end{center}
\end{cor}

\begin{proof}
$X = \Setabs{n \in N}{\phi(n)}$ لو، تے ہُن
\olref{thm:dedinfiniteinduction} ورتو۔
\end{proof}

ایس نتیجے وچ اسی فارمولے دے ”پیرامیٹر رکھن“ دی گل کیتی۔ لگ بھگ ایہدا
مطلب ایہ اے کہ کوئی وی شےواں $c_1, \ldots, c_k$ لئی اسی
$\phi(x, c_1, \ldots, c_k)$ نال کم کر سکدے آں۔ ہور ٹھیک طور تے اسی
”پیرامیٹر“ دا ذکر کیتے بنا نتیجہ ایویں بیان کر سکدے آں۔ ہر فارمولے
$\phi(x, v_1, \ldots, v_k)$ لئی، جیہدے سارے آزاد متغیر وکھائے گئے
نیں، ساڈے کول ایہ اے:
	\begin{align*}
			\forall v_1 \ldots \forall v_k((&\phi(o, v_1,\ldots, v_k) \land {}\\
			&	(\forall x \in N)(\phi(x,v_1, \ldots, v_k) \lif \phi(s(x), v_1,\ldots, v_k))) \lif {}\\
			&\hspace{3em} (\forall x \in N)\phi(x, v_1,\ldots, v_k))
	\end{align*}
صاف اے کہ ”پیرامیٹر رکھن“ دی گل لکھت نوں پڑھنا بہت سوکھا بنا سکدی اے۔
(\olref[sth][][]{part} وچ اسی ایہ طریقہ بہتی واری ورتاں گے۔)

ڈیڈیکائنڈ الجبریاں ول واپس آؤ: کوئی وی ڈیڈیکائنڈ الجبرا دِتا ہووے،
اسی جمع، ضرب تے قوت چڑھاؤن دے عام حسابی فنکشن وی تعریف کر سکدے آں۔
پر ایہ معمولی کم نہیں، تے ایہدے وچ \emph{بازگشتی تعریف} دی ترکیب
ورتدی اے۔ اسی ایہ ترکیب بہت بعد وچ، تے بہت ہور عام سیاق وچ، متعارف
کرا کے درست ثابت کراں گے۔ (دلچسپی رکھن والے
\olref[sth][ord-arithmetic][]{chap} توں بعد ایتھے مڑ آ سکدے نیں، یا
شاید \citealt[pp.~95--8]{Potter2004} ورگا کوئی متبادل بیان پڑھ لَین۔)
پر جتھے $N, s, o$ رل کے اک ڈیڈیکائنڈ الجبرا بناؤن، آخرکار اسی ایہ
مقرر کر سکاں گے:
\begin{align*}
	{a} + {o} &= {a} & & & {a} \times {o} &= {o} & & & {a}^{o} &= s(o)\\
{a} + {s}({b}) &= {s}({a}+{b}) &&& {a} \times {s}({b}) &= ({a}\times {b}) + {a}  & & & {a}^{{s}({b})} &= {a}^{b} \times {a}
\end{align*}
تے وکھا سکاں گے کہ ایہ اوویں ورتاؤ کردے نیں جیویں اسی امید کردے آں۔

\end{document}
